\chapter{Phân tích kết quả}
\section{Độ chính xác và mức bất định}
Trên tập kiểm thử C, SVM có tỷ lệ đúng toàn biển cao nhất trong các phương án được đối chiếu. So với SVM, nhánh Gaussian và HMM giảm 2,22 điểm phần trăm tỷ lệ đúng toàn biển. Khoảng tin cậy 95\% của chênh lệch là từ âm 5,56 đến dương 0,56 điểm phần trăm, tính bằng bootstrap ghép cặp 10\,000 lần theo ảnh. Khoảng này đi qua 0 nên chưa đủ cơ sở khẳng định sự khác biệt tỷ lệ đúng toàn biển trong tổng thể.

Ngược lại, CER của Gaussian và HMM cao hơn 6,15 điểm phần trăm, với khoảng tin cậy từ 3,82 đến 8,69 điểm phần trăm. Trong phép lấy mẫu lại trên tập ảnh hiện tại, mức tăng lỗi chuỗi có xu hướng nhất quán hơn chênh lệch số biển đúng. Nghiên cứu chưa định trước biên chấp nhận sai số để thực hiện kiểm định không thua kém, nên không kết luận hai phương án tương đương chất lượng.

Khoảng tin cậy tốc độ tương đối của Gaussian và HMM so với SVM là từ 1,307 đến 1,343 lần. Khoảng này phản ánh lấy mẫu lại theo ảnh trong cùng điều kiện đo, không mô tả biến thiên giữa các máy, camera hoặc phiên vận hành. Những nhận xét thống kê trên vẫn chịu giới hạn do tập kiểm thử đã được quan sát trong lịch sử.

\section{Sai số chuỗi và phân nhóm lỗi}
\begin{table}[H]\centering
\caption{Phân rã khoảng cách chỉnh sửa trên 1\,464 ký tự nhãn}
\begin{tabularx}{\textwidth}{Y r r r r}
\toprule
Phương án & Thay thế & Xóa & Chèn & Tổng\\
\midrule
\tablerows{error_results}
\bottomrule
\end{tabularx}
\end{table}

Gaussian và HMM có ít phép xóa hơn SVM, nhưng số phép thay thế tăng từ 159 lên 341. Kết quả cho thấy việc tạo chuỗi có độ dài gần nhãn hơn chưa bảo đảm nội dung đúng hơn. Cùng đúng 70 biển, nhánh SVM giảm chiều có 444 phép chỉnh sửa, còn nhánh Gaussian có 511 phép. Đây là lý do cần dùng đồng thời tỷ lệ đúng toàn biển và CER.

Nhóm phân loại các ca sai theo thứ tự chẩn đoán: trước hết xét vùng được chọn có IoU dưới 0,5; trong phần còn lại xét độ dài chuỗi khác nhãn; cuối cùng là vùng đạt ngưỡng, cùng độ dài nhưng nội dung sai. Theo cách này, 106 ca sai của SVM gồm 46 ca vùng không đạt, 48 ca khác độ dài và 12 ca cùng độ dài nhưng sai nội dung. Với Gaussian và HMM, 110 ca sai được chia lần lượt thành 52, 44 và 14 ca.

Đây là các nhóm dấu hiệu quan sát được, chưa phải phân rã nguyên nhân độc lập. Chẳng hạn, vùng đủ IoU vẫn có thể cắt mất nét ở mép; chuỗi thiếu ký tự có thể do phân đoạn hoặc do chọn giả thuyết khác. Không quy toàn bộ lỗi khác độ dài cho riêng bộ phân đoạn khi chưa có nhãn ký tự thủ công và kiểm tra trung gian.

\section{Các trường hợp điển hình}
\fig[1.0]{vi_du_loi.png}{Ví dụ đối chiếu kết quả. Khung xanh lá là nhãn vị trí, xanh lam là vùng SVM chọn và nâu là vùng Gaussian kết hợp HMM chọn; các khung có thể chồng nhau.}{fig:errors}

Với biển 36A43014, cả hai phương án chọn vùng phía trên xe thay vì biển ở cản trước. Đây là bằng chứng trực tiếp về lỗi lựa chọn vùng. Khi đầu vào OCR là một cấu trúc ngoài biển số, bộ phân loại vẫn có thể gán lớp và xuất chuỗi có vẻ giống ký tự. Cải thiện OCR riêng lẻ không giải quyết được trường hợp này nếu vùng sai vẫn được ưu tiên.

Với biển 36A40043, hai phương án chọn cùng khung vùng, nhưng Gaussian và HMM trả thêm ký tự 2 ở cuối. Sự khác nhau phù hợp với khả năng thay đổi giả thuyết phân đoạn hoặc điểm chất lượng OCR. Từ chuỗi và khung cuối chưa thể khẳng định một nguyên nhân duy nhất; cần kiểm tra từng vùng ký tự và điểm của các giả thuyết để xác định thành phần được đọc thừa.

Với biển 89G103316, SVM trả đúng nhãn còn Gaussian và HMM trả 89G13316, thiếu ký tự 0. Ảnh có bố cục hai hàng, hình chữ bị biến dạng theo góc nhìn và vùng biển có viền. Trường hợp này cho thấy lỗi thiếu ký tự vẫn tồn tại sau hiệu chỉnh và giải mã. HMM hiện tại chỉ chọn trạng thái cho các quan sát đã có, không có cơ chế tự tạo lại một ảnh ký tự bị thiếu.

Với biển 35G102993, SVM trả 3G12993 trong khi Gaussian và HMM trả đúng chuỗi. Hai phương án không khác khung vùng cuối. Ví dụ thể hiện mô hình quan sát và giải mã có thể giúp lựa chọn một giả thuyết nhận dạng tốt hơn, đồng thời không chứng minh Gaussian ưu thế trên toàn tập. Việc trình bày cả ca cải thiện và ca suy giảm giúp tránh lựa chọn minh họa chỉ có lợi cho một phương án.

\section{Phân tích ma trận nhầm lẫn}
Trên 434 ký tự chẩn đoán, SVM sai 38 mẫu còn Gaussian sai 46 mẫu. Với SVM, các cặp có hai lỗi gồm 7 thành 9, 6 thành 1 và 2 thành 3. Ở mô hình Gaussian, có bốn mẫu 1 bị dự đoán thành 7, hai mẫu 1 thành T và hai mẫu D thành 0. Các nhóm nét gần nhau này là vị trí cần xem lại biểu diễn hình dạng và số lượng mẫu.

Macro F1 của Gaussian là 77,99\%, thấp hơn accuracy 89,40\%. Chênh lệch phù hợp với việc có lớp hiếm và hiệu quả không đồng đều theo lớp. Một số lớp chỉ có một mẫu trong tập chẩn đoán; thay đổi một dự đoán có thể làm F1 lớp biến động mạnh. Những hàng không có nhãn thật trong ma trận không cung cấp bằng chứng về khả năng nhận dạng của lớp đó.

Các ma trận đo mô hình quan sát trước HMM. Không thể từ chúng kết luận HMM đã sửa hoặc gây ra từng lỗi. Muốn phân tích giải mã cần đối chiếu chuỗi thô và chuỗi sau HMM trên cùng giả thuyết, đồng thời theo dõi khả năng thay đổi thứ hạng vùng và phân đoạn.

\section{Chi phí tính toán và ưu, nhược điểm}
\fig{thoi_gian_cong_doan.png}{Thời gian từng công đoạn của các phương án A, B, C như Hình \ref{fig:mainresults}. Mỗi công đoạn lấy trung vị riêng theo ảnh nên tổng cột có thể khác thời gian toàn hệ thống.}{fig:stages}

Thời gian mô hình quan sát giảm từ khoảng 47,38 ms/ảnh ở SVM xuống 4,76 ms/ảnh ở Gaussian, bao gồm phép chiếu. Tuy nhiên, phát hiện và phân đoạn của nhánh Gaussian vẫn chiếm khoảng 113,81 ms/ảnh. Vì thế mức tăng tốc toàn hệ thống nhỏ hơn nhiều so với chênh lệch đo lô ký tự đã có sẵn HOG. Tối ưu tiếp theo cần tập trung cả phần sinh vùng và số giả thuyết phân đoạn.

SVM có ưu điểm là tỷ lệ đúng toàn biển và CER tốt nhất trong đối sánh hiện tại. Hạn chế là chi phí tính điểm quan sát cao hơn và không trực tiếp khai thác quan hệ chuỗi khi chỉ phân loại từng ký tự. PCA và FLD giảm kích thước biểu diễn, nhưng có nguy cơ làm mất hướng phân biệt cần thiết. Gaussian thuận lợi về tốc độ và cách diễn giải thống kê, song giả định một phân bố trên mỗi lớp có thể quá đơn giản với dữ liệu nhiều kiểu nét.

HMM cho phép kết hợp điểm ảnh với quan hệ ký tự và cấu trúc biển. Hạn chế là phụ thuộc mẫu cấu trúc, chất lượng xác suất quan sát và dữ liệu chuỗi train. Khi điểm ảnh sai hoặc có chuyển trạng thái thiên lệch, ràng buộc chuỗi có thể làm thay đổi một ký tự đúng thành sai. Trên nhánh SVM giảm chiều, kết quả bốn nhánh cho thấy HMM chưa mang lại lợi ích tổng thể.

Với mục tiêu ưu tiên độ chính xác hiện tại, nhóm chọn SVM cho chức năng mặc định. Nhánh PCA, FLD, Gaussian và HMM được giữ để chứng minh việc vận dụng kiến thức trong môn và nghiên cứu đánh đổi tốc độ với sai số. Nếu ứng dụng đòi hỏi phản hồi nhanh hơn, cần xác định trước mức lỗi chấp nhận được và kiểm tra trên tập độc lập trước khi đổi mô hình sử dụng.
